\chapter{Containers -- Docker}
\label{cap:implantacao}

\section{Imagens}

Todos os serviços usam o mesmo \textit{Dockerfile}, em duas fases e com imagens oficiais da
Microsoft: a compilação é feita na imagem do SDK do .NET 10 e o resultado é copiado para a imagem
\texttt{aspnet:10.0}, mais pequena e sem ferramentas de compilação. O processo corre com um
utilizador sem privilégios de administrador. O PostgreSQL, o RabbitMQ e o Jaeger usam também as
imagens oficiais.

\section{Ambiente com Docker Compose}

O ficheiro \texttt{deploy/docker-compose.yml} arranca o ambiente completo com um só comando:

\begin{verbatim}
docker compose --env-file .env -f deploy/docker-compose.yml up -d --build
\end{verbatim}

\begin{table}[H]
\centering
\small
\renewcommand{\arraystretch}{1.2}
\begin{tabularx}{\textwidth}{>{\raggedright\arraybackslash}p{0.28\textwidth} X}
\toprule
\textbf{Container} & \textbf{Função} \\
\midrule
\texttt{gateway} & API Gateway; único com porta publicada (8080) \\
\texttt{auth-api}, \texttt{restaurants-api}, \texttt{reservations-api}, \texttt{waitlist-api}, \texttt{notifications-api}
& Microsserviços, apenas na rede interna \\
\texttt{postgres}
& PostgreSQL 17; um \textit{script} de arranque cria as cinco bases de dados e um utilizador só de leitura para exportação de dados \\
\texttt{rabbitmq}
& \textit{Broker} de mensagens \\
\texttt{jaeger}
& Rastreio distribuído; interface apenas acessível na própria máquina \\
\bottomrule
\end{tabularx}
\caption{Containers do ambiente}
\label{tab:containeres}
\end{table}

Os serviços só arrancam depois de o PostgreSQL e o RabbitMQ estarem prontos, e cada serviço aplica as suas \textit{migrations} ao arrancar. As credenciais e a
chave do JWT vêm de um ficheiro \texttt{.env} que nunca é guardado no repositório; o repositório
inclui apenas um \texttt{.env.example}. Em ambiente de desenvolvimento, os serviços criam os dados de
exemplo do protótipo (três restaurantes em Braga e as contas de cliente, gestora e administrador).
